\documentclass[11pt]{article}

\usepackage[a4paper,margin=1in]{geometry}
\usepackage{amsmath,amssymb,amsthm}
\usepackage{booktabs}
\usepackage{microtype}
\usepackage{url}

\newtheorem{theorem}{Theorem}
\newtheorem{lemma}[theorem]{Lemma}
\newtheorem{corollary}[theorem]{Corollary}
\newtheorem{remark}[theorem]{Remark}

\newcommand{\dd}{\,\mathrm d}

\title{Conditional Implications Among the Unresolved Atlas Rows}
\author{}
\date{}

\begin{document}

\maketitle

\section{Result}

For a finite graph $G$ and $0\leq p<1$, put
\[
 \Phi_G(p)=(1-p)^{v(G)}
 \chi_G\!\left(\frac1{1-p}\right),
\]
with its polynomial continuation at $p=1$.  Let $F_{43}$ denote the
house graph (NetworkX Atlas $43$, graph6 \texttt{Dlc}) and let $F_{196}$
denote NetworkX Atlas $196$ (graph6 \texttt{ER\string~g}).

The following is a genuine implication between two rows that are not yet
completely classified.

\begin{theorem}[The blue arrow $43\longrightarrow196$]
\label{thm:arrow}
Assume that every graphon $V$ of edge density $z\in[1/2,1]$ satisfies
\[
 t(F_{43},V)\geq \Phi_{F_{43}}(z).
\]
Then every graphon $W$ of edge density $p\in[2/3,1]$ satisfies
\[
 t(F_{196},W)\geq \Phi_{F_{196}}(p).
\]
Thus a complete positive classification of Atlas $43$ would immediately
give a complete positive classification of Atlas $196$.
\end{theorem}

This is not merely a graph-containment observation.  Atlas $196$ is the
cone $K_1\vee F_{43}$, and the proof below transfers the assumed inequality
through normalized link graphons.  In particular, it applies to arbitrary
graphons, not only finite-step or regular graphons.

\section{The weighted link estimate}

Let $W$ be a graphon on a probability space $(\Omega,\mu)$, and write
\[
 p=\int_{\Omega^2}W(x,y)\dd\mu(x)\dd\mu(y),
 \qquad
 d(x)=\int_\Omega W(x,y)\dd\mu(y),
\]
and
\[
 \tau(x)=\int_{\Omega^2}
 W(x,y)W(x,z)W(y,z)\dd\mu(y)\dd\mu(z).
\]

\begin{lemma}[Degree-weighted rooted Goodman inequality]
\label{lem:weighted}
If $p>0$, then for every integer $k\geq0$,
\[
 \int_\Omega d(x)^k\tau(x)\dd\mu(x)
 \geq \frac{2p-1}{p}
       \int_\Omega d(x)^{k+2}\dd\mu(x).
\]
\end{lemma}

\begin{proof}
Put $M_j=\int d^j$, and define
\[
 I_k=\int d(x)^k\tau(x)\dd\mu(x),
 \qquad
 J_k=\int_{\Omega^2}d(x)^kW(x,y)d(y)
       \dd\mu(x)\dd\mu(y).
\]
The elementary inequality $ab\geq a+b-1$ on $[0,1]^2$, followed by
integration, gives
\[
 I_k\geq 2J_k-pM_k.
\]
Set $U=1-W$, $u=1-d$, and $q=1-p$.  If $T_U$ is the integral operator
with kernel $U$, then
\[
 T_Wd=d-q+T_Uu.
\]
Moreover, symmetry of $U$ gives
\[
\begin{split}
 2\left(\int d^kT_Uu-\int d^ku^2\right)
 ={}&\int_{\Omega^2}U(x,y)
       \bigl(u(y)-u(x)\bigr)
       \bigl(d(x)^k-d(y)^k\bigr)\dd\mu^2\\
 \geq{}&0,
\end{split}
\]
because $u=1-d$.  It follows that
\[
 J_k\geq M_{k+2}-M_{k+1}+pM_k
\]
and hence
\[
 I_k\geq2M_{k+2}-2M_{k+1}+pM_k.
\]
Finally,
\[
 2M_{k+2}-2M_{k+1}+pM_k
 -\frac{2p-1}{p}M_{k+2}
 =\frac1p\int d(x)^k(d(x)-p)^2\dd\mu(x)\geq0.
\]
\end{proof}

For $d(x)>0$, let $W_x$ be the same kernel $W$ on the probability measure
\[
 \dd\nu_x(y)=\frac{W(x,y)}{d(x)}\dd\mu(y),
\]
and let $z_x=t(K_2,W_x)=\tau(x)/d(x)^2$.  Put $z_x=0$ when $d(x)=0$.
All link-integral contributions on $\{x:d(x)=0\}$ are taken to be zero.
Lemma~\ref{lem:weighted}, with $k=3$, says
\begin{equation}
 \label{eq:link-average}
 \int d(x)^5z_x\dd\mu(x)
 \geq \left(2-\frac1p\right)\int d(x)^5\dd\mu(x).
\end{equation}

\section{Proof of the implication}

The house target is
\[
 \phi(z):=\Phi_{F_{43}}(z)
 =z(2z-1)(3z^2-3z+1)
 =6z^4-9z^3+5z^2-z.
\]
Its derivatives are
\[
 \phi'(z)=24z^3-27z^2+10z-1,
 \qquad
 \phi''(z)=2(3z-1)(12z-5).
\]
Thus $\phi''>0$ on $[1/2,1]$, and
$\phi'(1/2)=1/4$, so $\phi'$ is positive there.  Also
$\phi\geq0$ on this interval and $\phi(1/2)=0$.

\begin{proof}[Proof of Theorem~\ref{thm:arrow}]
The endpoint $p=1$ is immediate, so suppose $2/3\leq p<1$ and put
\[
 c=2-\frac1p\in[1/2,1].
\]
Let
\[
 \ell_c(z)=\phi(c)+\phi'(c)(z-c)
\]
be the tangent to $\phi$ at $c$.  We claim that every graphon $V$ of
arbitrary edge density $z\in[0,1]$ satisfies
\begin{equation}
 \label{eq:tangent-minorant}
 t(F_{43},V)\geq\ell_c(z).
\end{equation}
If $z\geq1/2$, this follows from the conditional hypothesis and convexity:
$t(F_{43},V)\geq\phi(z)\geq\ell_c(z)$.  If $z<1/2$, then
$\phi'(c)\geq0$ and convexity gives
\[
 \ell_c(z)\leq\ell_c(1/2)\leq\phi(1/2)=0
 \leq t(F_{43},V).
\]

Atlas $196$ is $K_1\vee F_{43}$.  Integrating first over the five
neighbours of its universal vertex gives the exact link identity
\[
 t(F_{196},W)
 =\int_\Omega d(x)^5t(F_{43},W_x)\dd\mu(x).
\]
Apply \eqref{eq:tangent-minorant} to every $W_x$.  Using
\eqref{eq:link-average} and $\phi'(c)\geq0$, we obtain
\[
\begin{split}
 t(F_{196},W)
 &\geq \phi(c)\int d^5
 +\phi'(c)\left(\int d^5z_x-c\int d^5\right)\\
 &\geq \phi(c)\int d^5.
\end{split}
\]
Jensen gives $\int d^5\geq p^5$, and $\phi(c)\geq0$, hence
\[
 t(F_{196},W)\geq p^5\phi(c).
\]

It remains only to identify the target.  With $q=1-p$, one has
$1-c=q/p$.  Also
\[
 \chi_{K_1\vee F_{43}}(x)=x\chi_{F_{43}}(x-1).
\]
Consequently
\[
 p^5\phi(c)
 =q^5\chi_{F_{43}}\!\left(\frac pq\right)
 =\Phi_{K_1\vee F_{43}}(p)
 =\Phi_{F_{196}}(p),
\]
which proves the theorem.
\end{proof}

For reference, the exact chromatic polynomials are
\[
 \chi_{F_{43}}(x)=x(x-1)(x-2)(x^2-3x+3)
\]
and
\[
 \chi_{F_{196}}(x)
 =x(x-1)(x-2)(x-3)(x^2-5x+7).
\]
The latter target is therefore
\[
 \Phi_{F_{196}}(p)
 =p(2p-1)(3p-2)(7p^2-9p+3).
\]

One explicit Atlas identification is as follows.  In the NetworkX
labelling of Atlas $43$, the edge list is
\[
 01,\ 03,\ 04,\ 12,\ 23,\ 34.
\]
Add a universal vertex $5$.  The map
\[
 0\mapsto2,\quad 1\mapsto3,\quad 2\mapsto1,\quad
 3\mapsto5,\quad 4\mapsto0,\quad 5\mapsto4
\]
carries the resulting cone to the NetworkX edge list of Atlas $196$.

\section{Quantitative propagation of partial progress}

The same argument explains the already accepted partial range for Atlas
$196$.  Suppose the house inequality is known only for $z\geq a$, where
$a\in[1/2,5/6)$, while $\phi,\phi',\phi''\geq0$ on $[a,1]$.  Define
\[
 L_a(c)=\phi(c)+(a-c)\phi'(c).
\]
If $c_a\in[a,1]$ is the least point for which $L_a(c_a)\leq0$, then the
same tangent proof gives
\[
 \boxed{
  \text{Atlas }43\text{ proved for }z\geq a
  \quad\Longrightarrow\quad
  \text{Atlas }196\text{ proved for }
  p\geq\frac1{2-c_a}.}
\]
For the current house threshold
\[
 a=0.7085121949942608\ldots,
\]
this gives
\[
 c_a=0.8412137423978822\ldots,
 \qquad
 \frac1{2-c_a}=0.8629719186257050\ldots,
\]
which is exactly the green partial range recorded for Atlas $196$.
When $a=1/2$, one has $L_a(a)=\phi(a)=0$, so $c_a=a$ and the transferred
threshold is $1/(2-a)=2/3$, proving Theorem~\ref{thm:arrow} on the entire
required interval.

\section{Audit of other possible arrows}

\subsection{Three exact range-transfer arrows}

There are also three exact implications obtained from spanning containment.
Unlike Theorem~\ref{thm:arrow}, each implication covers only a proper initial
subinterval of the target row's required range.

\begin{theorem}[Spanning-supergraph range transfers]
\label{thm:range-arrows}
The following statements hold for arbitrary graphons.
\begin{enumerate}
\item If Atlas $188$ satisfies its chromatic-polynomial bound on its full
required interval, then Atlas $171$ satisfies its bound for
\[
 \frac12\leq p\leq\frac23.
\]
\item Under the same hypothesis on Atlas $188$, Atlas $153$ satisfies its
bound for
\[
 \frac12\leq p\leq\frac35.
\]
\item If Atlas $199$ satisfies its chromatic-polynomial bound on its full
required interval, then Atlas $181$ satisfies its bound for
\[
 \frac23\leq p\leq\frac34.
\]
\end{enumerate}
None of these statements claims anything about the remainder of the target row's
required interval.
\end{theorem}

\begin{proof}
In the NetworkX labellings, the vertex map
\[
 \sigma:(0,1,2,3,4,5)\longmapsto(0,5,4,3,1,2)
\]
embeds Atlas $171$ as the spanning subgraph of Atlas $188$ obtained by
deleting the single edge $15$.  Consequently, since $0\leq W\leq1$,
\begin{equation}
 t(F_{171},W)\geq t(F_{188},W).                 \label{eq:171-188-density}
\end{equation}
Their targets satisfy the exact identity
\begin{equation}
 \Phi_{188}(p)-\Phi_{171}(p)
 =p(p-1)(2p-1)^2(3p-2).                        \label{eq:171-188-target}
\end{equation}
For $1/2\leq p\leq2/3$, the two factors $p-1$ and $3p-2$ are both
nonpositive and every other factor is nonnegative.  Thus
$\Phi_{188}(p)\geq\Phi_{171}(p)$ there.  The assumed Atlas $188$ bound,
\eqref{eq:171-188-density}, and \eqref{eq:171-188-target} give
\[
 t(F_{171},W)\geq t(F_{188},W)
 \geq\Phi_{188}(p)\geq\Phi_{171}(p).
\]

The vertex map
\[
 \rho:(0,1,2,3,4,5)\longmapsto(0,1,3,4,5,2)
\]
embeds Atlas $153$ as the spanning subgraph of Atlas $188$ obtained by
deleting the two edges $15$ and $24$.  Thus
\begin{equation}
 t(F_{153},W)\geq t(F_{188},W).                 \label{eq:153-188-density}
\end{equation}
The exact target comparison is
\begin{equation}
 \Phi_{188}(p)-\Phi_{153}(p)
 =p(p-1)(2p-1)^2(5p-3).                        \label{eq:153-188-target}
\end{equation}
For $1/2\leq p\leq3/5$, the factors $p-1$ and $5p-3$ are both
nonpositive and all other factors are nonnegative.  Hence the same
three-step comparison gives the second assertion.  This is an exact but
presently redundant transfer: the unconditional Atlas $153$ range already
extends beyond $3/5$ by a direct argument.

Similarly, the map
\[
 \eta:(0,1,2,3,4,5)\longmapsto(1,3,2,4,5,0)
\]
embeds Atlas $181$ as the spanning subgraph of Atlas $199$ obtained by
deleting the single edge $03$.  Hence
\begin{equation}
 t(F_{181},W)\geq t(F_{199},W).                 \label{eq:181-199-density}
\end{equation}
Here the target difference factors as
\begin{equation}
 \Phi_{199}(p)-\Phi_{181}(p)
 =p(p-1)(2p-1)(3p-2)(4p-3).                    \label{eq:181-199-target}
\end{equation}
On $2/3\leq p\leq3/4$, the factors $p-1$ and $4p-3$ are both
nonpositive, while the remaining factors are nonnegative.  Therefore
$\Phi_{199}(p)\geq\Phi_{181}(p)$, and the same three-step comparison proves
the third assertion.
\end{proof}

\begin{remark}[Why these arrows are range-qualified]
For $p>2/3$, the right-hand side of \eqref{eq:171-188-target} is negative;
for $p>3/5$, the right-hand side of \eqref{eq:153-188-target} is negative;
for $p>3/4$, the right-hand side of \eqref{eq:181-199-target} is negative.
Thus plain spanning containment reverses the needed target comparison beyond
the displayed endpoints.  These are rigorous partial dependencies, not
complete classification arrows.
\end{remark}

The exact audit script exhausts all $108$ ordered spanning-subgraph
containments among the twenty unresolved rows.  It factors every target gap
over the rationals; every irreducible factor is linear or a quadratic with
negative discriminant, so rational sign-cell tests are exact.  After
intersecting the target's required interval with the interval on which the
source theorem is assumed, the three implications in
Theorem~\ref{thm:range-arrows} are the only positive-length loci.  Four raw
target-gap loci ($194\to153$, $194\to171$, $196\to153$, and $196\to171$)
lie wholly below the sources' required threshold $p=2/3$ and therefore do
not constitute conditional implications.

\subsection{Closure-operation audit}

The current unresolved set consists of the six partial rows
\[
 43,153,171,174,188,196
\]
and the fourteen open rows
\[
 118,122,124,130,147,151,157,168,169,181,185,194,199,203.
\]
An exact NetworkX audit of the accepted closure operations gives the
following result.

\begin{center}
\begin{tabular}{@{}p{0.25\linewidth}p{0.50\linewidth}p{0.18\linewidth}@{}}
\toprule
Operation & order restriction for a target on at most $6$ vertices
          & unresolved-base arrow\\
\midrule
add one universal vertex & base has at most $5$ vertices
                         & $43\longrightarrow196$\\
full whiskering & base has at most $3$ vertices & none\\
vertex self-amalgam & every nontrivial unresolved base is too large & none\\
edge self-amalgam & every nontrivial unresolved base is too large & none\\
disjoint union & no applicable connected unresolved target & none\\
\bottomrule
\end{tabular}
\end{center}

There are three especially tempting non-arrows.  Atlas $118$, $122$, and
$124$ are obtained by attaching one leaf to the three vertex orbits of the
house, and all three have target $p\phi(p)$.  However, a one-leaf operation
is not covered by full whiskering.  If $f_r(x)$ is the house density rooted
at the attachment vertex $r$, the assumed house theorem controls only
\[
 \int f_r(x)\dd\mu(x),
\]
whereas the required leaf density is
\[
 \int d(x)f_r(x)\dd\mu(x).
\]
The unrooted assertion supplies no sign for the covariance term
\[
 \int (d(x)-p)f_r(x)\dd\mu(x).
\]
Therefore this note does \emph{not} assert any of the arrows
$43\to118$, $43\to122$, or $43\to124$.  Proving one of them requires a
new rooted or degree-biased house theorem, not just the present unrooted
Atlas $43$ statement.

The audit is reproduced by
\texttt{experiments/open\_conditional\_implication\_audit.py}.  Its purpose
is to prevent a visual subgraph relation from being mistaken for a valid
graphon closure theorem.

\section{Dependencies and scope}

The conditional proof uses only Tonelli--Fubini, the elementary weighted
rooted Goodman inequality proved above, convexity of the explicit house
target, Jensen's inequality, and the assumed complete Atlas $43$ theorem.
It uses neither flag algebras nor a finite-step reduction.  The theorem is
conditional: it does not change the present partial status of either row.

The general universal-vertex theorem is also proved in
\path{notes/clique_oddcycle_join.tex}; the quantitative delayed-range
version is proved in
\path{notes/atlas196_partial_house_cone_lift.tex}.  The proof here is
included so that the blue dependency arrow has a single, self-contained
reference.

\end{document}
